\chapter{Supplementary Materials}
\label{app:supplementary}

\section{Synthetic Proposal Ground Truth}
\label{app:groundtruth}

Table~\ref{tab:groundtruth} presents the ground truth annotations for all 20 synthetic proposals.

\begin{longtable}{llp{5cm}l}
\caption{Ground truth annotations for synthetic evaluation proposals.}
\label{tab:groundtruth} \\
\toprule
\textbf{ID} & \textbf{Risk} & \textbf{Domain} & \textbf{Key Rights} \\
\midrule
\endfirsthead
\toprule
\textbf{ID} & \textbf{Risk} & \textbf{Domain} & \textbf{Key Rights} \\
\midrule
\endhead
P01 & High & Facial recognition in education & Art. 7, 8, 21 \\
P02 & High & Predictive policing with personal data & Art. 6, 8, 21, 47 \\
P03 & High & Healthcare diagnosis with patient records & Art. 2, 8, 21, 35 \\
P04 & High & Automated hiring with demographic data & Art. 15, 21, 23 \\
P05 & High & Social credit scoring system & Art. 1, 7, 8, 21 \\
P06 & High & Child behaviour monitoring AI & Art. 7, 8, 24 \\
P07 & High & Autonomous weapons targeting & Art. 2, 6 \\
P08 & High & Immigration decision automation & Art. 18, 19, 21, 47 \\
P09 & High & Mental health prediction from social media & Art. 7, 8, 35 \\
P10 & Medium & Sentiment analysis on public tweets & Art. 8, 11 \\
P11 & Medium & Traffic flow optimisation & Art. 6, 8 \\
P12 & Medium & Agricultural crop prediction & Art. 8 \\
P13 & Medium & Customer service chatbot & Art. 8, 38 \\
P14 & Medium & Energy consumption forecasting & Art. 8 \\
P15 & Medium & Language translation tool & Art. 11, 22 \\
P16 & Medium & Plagiarism detection in education & Art. 8, 14 \\
P17 & Low & Weather pattern analysis & Art. 8 \\
P18 & Low & Historical document digitisation & Art. 22 \\
P19 & Low & Mathematical theorem proving & -- \\
P20 & Low & Open-source code linting tool & -- \\
\bottomrule
\end{longtable}


\section{SPARQL Competency Queries}
\label{app:sparql}

\subsection{CQ1: Incidents Impacting Non-Discrimination}

\begin{lstlisting}[language=SPARQL,caption={CQ1: Retrieve incidents impacting Article 21.}]
PREFIX aief: <https://w3id.org/aief/>
SELECT ?incident ?label
WHERE {
  ?incident a aief:AIIncident ;
            aief:impactsRight aief:Article21NonDiscrimination ;
            rdfs:label ?label .
}
ORDER BY ?incident
\end{lstlisting}

\subsection{CQ2: Requirements Mapped to Data Protection}

\begin{lstlisting}[language=SPARQL,caption={CQ2: Retrieve requirements mapped to Article 8.}]
PREFIX aief: <https://w3id.org/aief/>
SELECT ?req ?label ?framework
WHERE {
  ?req a aief:GovernanceRequirement ;
       aief:mapsToRight aief:Article8DataProtection ;
       rdfs:label ?label ;
       dct:source ?framework .
}
ORDER BY ?req
\end{lstlisting}

\subsection{CQ3: Charter Rights for a Given Incident}

\begin{lstlisting}[language=SPARQL,caption={CQ3: Retrieve rights impacted by AIAAIC-001.}]
PREFIX aief: <https://w3id.org/aief/>
SELECT ?right ?label
WHERE {
  aief:AIAAIC-001 aief:impactsRight ?right .
  ?right rdfs:label ?label .
}
\end{lstlisting}

\subsection{CQ4: Incidents Supporting a Requirement}

\begin{lstlisting}[language=SPARQL,caption={CQ4: Retrieve incidents supporting R085.}]
PREFIX aief: <https://w3id.org/aief/>
SELECT ?incident ?label
WHERE {
  ?incident aief:supportsRequirement aief:R085 ;
            rdfs:label ?label .
}
\end{lstlisting}


\section{Ontology Class Hierarchy (Excerpt)}
\label{app:ontology}

The following presents the top-level class hierarchy of the AIEF ontology:

\begin{verbatim}
owl:Thing
  +-- aief:ResearchProposal
  +-- aief:GovernanceRequirement
  |     +-- aief:REAMSRequirement
  |     +-- aief:EUAIActRequirement
  |     +-- aief:HorizonEuropeRequirement
  |     +-- aief:ACMRequirement
  +-- aief:DeonticModality
  |     +-- aief:Obligation  [rdfs:subClassOf odrl:Obligation]
  |     +-- aief:Permission  [rdfs:subClassOf odrl:Permission]
  |     +-- aief:Prohibition [rdfs:subClassOf odrl:Prohibition]
  +-- aief:FundamentalRight  [owl:equivalentClass dpv-rights]
  +-- aief:AIIncident
  +-- aief:RiskLevel
  |     +-- aief:HighRisk
  |     +-- aief:MediumRisk
  |     +-- aief:LowRisk
  +-- aief:Mitigation
  +-- aief:AICapability
  +-- aief:Harm
  +-- aief:GovernanceFramework
\end{verbatim}


\section{Expert Review Scoring Sheet}
\label{app:scoring}

The scoring sheet prepared for expert review covers 8 proposals (P01, P06, P08, P13, P20, R01, R02, R03). For each proposal, reviewers are asked to assess the system's output on four dimensions using a 5-point Likert scale:

\begin{enumerate}
\item \textbf{Risk level correctness:} Is the assigned risk level (High/Medium/Low) appropriate for this proposal? (1 = Completely incorrect, 5 = Completely correct)
\item \textbf{Requirement relevance:} Are the cited requirements relevant to the proposal? (1 = None relevant, 5 = All relevant)
\item \textbf{Completeness:} Does the assessment cover all significant ethics issues? (1 = Major gaps, 5 = Comprehensive)
\item \textbf{Rights mapping accuracy:} Are the identified Charter rights correctly linked to the proposal's risks? (1 = Completely incorrect, 5 = Completely correct)
\end{enumerate}

This scoring instrument was designed to enable computation of Spearman $\rho$ for inter-rater agreement if multiple expert responses are received.
